\begin{table}[t]
\centering
\caption{Mode against size. One fixed pool of 2,500 PBMC14K cells, drawn afresh at each of three seeds, run through both of MICA's pipelines and the baseline under the same rule. The data, the size, the subsamples and the labels are identical across the three rows. The pipelines differ, and so do the cluster counts their grids reach: graph mode cannot make fewer than 4 or 5 clusters here, and Section~\ref{sec:modesize} shows that this, not the mode, accounts for the difference. Tied settings are those within 0.01 of the best. The span is the range of cluster counts among them, reported as the median over the three pools with the observed range. The last two columns say how often the two readings of the rule chose differently and by how much in ARI.}
\label{tab:modesize}
\footnotesize
\setlength{\tabcolsep}{4pt}
\resizebox{\columnwidth}{!}{%
\begin{tabular}{lcccc}
\toprule
Pipeline & Share tied & Span & Readings differ & Median gap \\
\midrule
MICA, small-data mode & 0.26 & 4 (2\textrm{-}5) & 3 of 3 & 0.319 \\
MICA, graph mode & 0.16 & 0 (0\textrm{-}1) & 0 of 3 & 0.000 \\
Leiden, as a control & 0.15 & 1 (1\textrm{-}1) & 1 of 3 & 0.000 \\
\bottomrule
\end{tabular}}
\end{table}
